\section{Exact fibers of a symbolic quotient}\label{sec:quotient-fibers}\label{sec:quotients}

Our next objective is a bijection between $A$ and each of its
nonwellorderable quotients. A local count of names is insufficient:
assignments made in two diagrams must agree whenever those diagrams
occur in one generic filter. We first classify local elements by the
subdiagrams on which they are available. We then construct bijections
which respect every restriction and every change of coordinates.

Throughout the quotient construction, the cardinal parameters are those
of Section~\ref{sec:forcing}, with the monoid $T$ supplied there by
Theorem~\ref{thm:monoid}. The forcing $\mathbb P$ remains the
ground set of diagrams, including when it is used after forcing with
$\mathbb R$. We use countable closure of $\mathbb R$, regularity of
$\kappa=\omega_1$, preservation of the ground cardinals $\mu,\mu^+$,
and the ground equality $\mu^{\aleph_0}=\mu$. Preservation of smaller
cardinals is unnecessary. In particular, $\mu$ always denotes the
chosen ground ordinal, not a cardinal recomputed from a new continuum.
All ground selections below are choices over specified sets in $V$.

Fix a hereditarily symmetric name $\dot E$ and a product condition
$(d_0,p)$ symmetrically forcing that $\dot E$ is an equivalence relation
on $A$ whose quotient is not wellorderable. Enlarge the diagram so that
$D_p$ supports $\dot E$ and contains $\supp_I(d_0)$. Choose $q$ with
$q_e=p$. Our bijection will have support $D_q$ and will be forced below
$(d_0,q)$. Write $G_q=\Fix(D_q)$. This group fixes $q,p,d_0$ and
$\dot E$ literally.

We write $F$ for an $\mathbb R$-generic filter containing $d_0$, and
work temporarily in $V[F]$. A product name $\tau$ determines an ordinary
$\mathbb P$-name $\tau[F]$ there. Thus an assertion such as
$r\forces\tau=\sigma$ in this section means ordinary $\mathbb P$-forcing
in $V[F]$, using these sections of the product names. The generic
notation abbreviates forcing arguments uniform below $d_0$; no claim
that a ground permutation acts on one fixed evaluated universe is used.

\subsection{Local names and their cardinal bound}

Recall the local name collection from
Equation~\eqref{eq:normalized-local-names}. A normalized subset name
means a subset of the ground set
\[
 \{(\dot A_j,(d,s)):j\in I,
                 \ (d,s)\in\mathbb R\times\mathbb P\}.
\]
For a diagram $r$, $\mathcal X(r)$ is the ground set of all such
names fixed literally by $\Fix(D_r)$. Each is hereditarily symmetric,
since its member names are the individual label names. Normalization by
the forcing relation retains any given support. When $r\le q$, define
in $V[F]$
\[
 \mathcal D_F(r)=
 \{\tau\in\mathcal X(r):
       r\forces\tau\text{ is an }\dot E\text{-class}\}/\mathord\sim_r,
 \qquad
 \tau\sim_r\sigma\ \Longleftrightarrow\ r\forces\tau=\sigma.
\]
These are quotients of ground sets, so they are sets. We sometimes omit
$F$ from the notation. Their elements are symbolic classes of names,
not yet the evaluated members of $A/E$ in a product extension.
Lemma~\ref{lem:local-decision} applies to the
class and equality tests. Consequently, for $r\le s\le q$, retaining
the same name defines an injection
\begin{equation}\label{eq:symbolic-inclusion}
 \iota_{s,r}:\mathcal D_F(s)\longrightarrow\mathcal D_F(r).
\end{equation}
Indeed, if a stronger diagram forced two local names equal, the decision
at $s$ would already be equality. These maps compose literally on
representatives. Every $j\in D_r$ has a normalized name, supported by
$D_r$, for its $E$-class: use the forcing relation to name the labels
$E$-related to $\dot A_j$. Below the present hypothesis it is an
eligible member of $\mathcal D_F(r)$.

\begin{lemma}[Local quotient bound]\label{lem:local-quotient-bound}
For every $F$ as above and every ground $r\le q$,
$|\mathcal D_F(r)|\le\mu$ in $V[F]$.
\end{lemma}
\begin{proof}
For an element choose a representative $\tau$ and a diagram $t\le r$
forcing $\tau=[\dot A_j]_E$ for some $j\in D_t$. The ordinary forcing
theorem permits deciding a representative label; extending the diagram
includes its index. Write $r=t_b$. By Cohen restriction in
Lemma~\ref{lem:local-decision}, there is $d\in F$, supported in $D_t$,
witnessing this $\mathbb P$-forcing assertion. Record the triple
\[
 (b,x,t^{-1}d),\qquad t(x)=j,
\]
where $t^{-1}d$ is the partial binary function on $T\times\kappa$
obtained by replacing each coordinate $t(y)$ with $y$. Its domain has
size less than $\kappa$. The ground set of possible records has size
at most $\mu$, since $|T|=\mu$, $\kappa\le\mu$, and
$\mu^{<\kappa}=\mu$. Its ground enumeration remains an enumeration in
$V[F]$.

Suppose two chosen representatives $\tau_1,\tau_2$ have the same
record, witnessed by $t_1,t_2,d_1,d_2,j_1,j_2$. The ground bijection
$t_1(y)\mapsto t_2(y)$ agrees with the identity on $D_r$: in both
diagrams $r=(t_i)_b$. It maps points outside $D_r$ to points outside
$D_r$. Extend it to a ground permutation $\pi$ of $I$ fixing $D_r$.
This is possible since the two ranges have size $\mu<|I|$ and their
complements have the same cardinality $|I|$. Then
\[
 \pi t_1=t_2,\qquad \pi d_1=d_2,\qquad \pi j_1=j_2,
 \qquad \pi\tau_1=\tau_1.
\]
Product symmetry makes $d_2$ witness at $t_2$ equality of both
$\tau_1$ and $\tau_2$ with $[\dot A_{j_2}]_E$. Since $d_2\in F$,
$t_2$ forces $\tau_1=\tau_2$ in $V[F]$. Decision at $r$ gives the
same equality there. Thus distinct elements have distinct records.
Choosing representatives and records uses Choice in the ordinary
extension $V[F]$, not in the symmetric model.
\end{proof}

\subsection{Availability and exact multiplicities}

For $r\le q$, write $q=r_a$. For $h\in\mathcal D_F(r)$ and
$j\in D_r$, respectively, define
\begin{align}
 U_h&=\{n\sqsubseteq a:
             h\in\ran(\iota_{r_n,r})\},\label{eq:class-availability}\\
 U_j&=\{n\sqsubseteq a:j\in D_{r_n}\}.
             \label{eq:label-availability}
\end{align}
An ideal here is nonempty, downward closed, and closed under finite
joins. All joins needed within $\downarrow a$ exist.

\begin{lemma}[Availability ideals]\label{lem:availability-ideals}
Each set in Equations~\eqref{eq:class-availability} and
\eqref{eq:label-availability} is an ideal of $\downarrow a$.
Every ideal of $\downarrow a$ in $V[F]$ belongs to $V$.
\end{lemma}
\begin{proof}
Both sets contain $1$ and are downward closed by inclusion of
subdiagrams. Suppose $n,n'\in U_h$ and put $k=n\vee n'$. The ranges
of $r_n,r_{n'}$ intersect exactly in $D_{r_k}$. Choose supported names
for their representations of $h$. They are forced equal by $r$.
Lemma~\ref{lem:support-descent}, with root $r_k$, supplies a
$D_{r_k}$-supported name equal to them. Normalize it without changing
support. Decision at $r_k$ shows it is an $E$-class there: its
negation would persist to the extension $r$. Thus $k\in U_h$.
For labels the identical argument is the literal equality of cone
intersections, without any name descent.

For the final assertion, apply Lemma~\ref{lem:mono-countable-ideal}
in $V[F]$. The ground monoid, order, heights and costs retain their
meanings there, so every nonempty ideal has a countable cofinal sequence.
It consists of ground elements of $T$, so countable closure of
$\mathbb R$ puts it in $V$. Its downward closure, computed with the
unchanged ground order, is the original ideal. The argument applies
to every ideal in $V[F]$, not only to ground ideals.
\end{proof}

For a nonempty ideal $U\subseteq\downarrow a$, put
\[
 B_{r,U}(F)=\{h\in\mathcal D_F(r):U_h=U\},\qquad
 L_{r,U}=\{j\in D_r:U_j=U\}.
\]
The second set is ground data; the first can depend on $F$.
These exact fibers, rather than the whole local sets $\mathcal D_F(r)$,
are the pieces on which a bijection can respect restriction: a label
and its assigned class must remain available at precisely the same
subdiagrams.

\begin{lemma}[Exact fibers]\label{lem:exact-fibers}
For every $r\le q$ and every nonempty ideal
$U\subseteq\downarrow a$, where $q=r_a$, both $B_{r,U}(F)$ and
$L_{r,U}$ have cardinality $\mu$ in $V[F]$. The label fiber already
has that cardinality in $V$.
\end{lemma}
\begin{proof}
The common root $p$ will separate the copies used in this proof.
We first find a class with no representation supported by $D_p$.
Copying it into diagrams whose unwanted intersections are exactly
$D_p$ will then rule out both coincidences between copies and
availability at forbidden subdiagrams, by support descent.

\smallskip\noindent\textit{A class unavailable at the root.}
Enumerate
$\mathcal X(p)$ in $V$ as $(\tau_\xi:\xi<\lambda)$. The canonical
name for the sequence of its evaluations is hereditarily symmetric
with support $D_p$. Explicitly, its entries are the canonical pairs
$(\check\xi,\tau_\xi)$, each placed under every product condition.
Every member of $\Fix(D_p)$ fixes each entry literally. The range
of this sequence is wellorderable inside $W$: map each value to its
least index of occurrence. If it contained every $E$-class, the same
map restricted to $A/E$ would wellorder that quotient. The hypothesis
therefore implies, ordinarily below $(d_0,p)$, that some $E$-class
lies outside the evaluations of all names in $\mathcal X(p)$.

Fix $F$ containing $d_0$. By the ordinary forcing theorem choose
$t\le p$ and $j\in D_t$ such that $t$ forces the class of $A_j$
outside that entire evaluation list. Choose a witnessing Cohen
condition $d\in F$ supported in $D_t$, using
Lemma~\ref{lem:local-decision}. Write $p=t_b$. Necessarily $b\ne1$,
since a class represented by a label in $D_p$ has a name in
$\mathcal X(p)$.

\smallskip\noindent\textit{Independent copies with active Cohen conditions.}
As $p=r_{ae}$, the probe property of Theorem~\ref{thm:monoid}
supplies $\mu$ elements $m$ with $mb=ae$ and the prescribed ideal
$U$. The diagrams $r_m$ are copies of $t$ over $p$, with pairwise
intersection exactly $D_p$. Choose the ground permutations taking
$t$ to these diagrams and fixing $D_p$, and transport $d$ and the
class name along them. Their nonroot coordinate sets are disjoint.
This is the role of the earlier extension $q_e=p$: the availability
pattern is prescribed among divisors of $a$, while the common root
lies at $ae$, strictly beyond $a$.

There are $\mu$ copies whose transported Cohen conditions belong to
$F$. To see this without any $\mu$-closure assertion, partition the
ground family into $\mu$ ground subfamilies of size $\mu$. In each
subfamily the conditions extending at least one transported copy of
$d$ form a dense set below $d\restriction D_p$. Indeed, a Cohen
condition uses countably many coordinate indices; among the pairwise
disjoint nonroot ranges there is one it misses, and on the root the
two conditions agree. Their union is a Cohen condition. Genericity
meets each of these ground dense sets. A hit in each disjoint
subfamily yields $\mu$ different hits.

\smallskip\noindent\textit{Distinct classes with the prescribed exact ideal.}
Each hit gives a $D_{r_m}$-supported class name forced at $r_m$
to lie outside the evaluations of $\mathcal X(p)$. The diagram
$r_m$ extends $p$ but need not extend $q$; we use its class name,
not a set $\mathcal D_F(r_m)$. As $r$ extends $r_m$, this name
represents a symbolic class in the defined set $\mathcal D_F(r)$.
Two such
classes cannot be equal there: support descent over their common
root $p$ would supply a name in $\mathcal X(p)$ for their common
value, contrary to the outside assertion. If $n\in U$, then
$n\sqsubseteq m$, so $r_n$ extends $r_m$ and the class is available
at $r_n$. If $n\sqsubseteq a$ and $n\notin U$, the exact probe
intersection gives
\[
 nT\cap mT=aeT,
 \qquad D_{r_n}\cap D_{r_m}=D_p.
\]
Availability at $r_n$ would again descend to a name in
$\mathcal X(p)$, a contradiction. Its exact availability ideal is
therefore $U$. This gives $\mu$ distinct elements of $B_{r,U}(F)$;
Lemma~\ref{lem:local-quotient-bound} bounds that fiber from above.

\smallskip\noindent\textit{The label fiber.}
For labels use the probe centers $r(m)$. Their availability ideals
are exactly $\downarrow m\cap\downarrow a=U$, and distinct probes
have distinct centers. This supplies a ground $\mu$-subset of
$L_{r,U}$; the upper bound follows from $|D_r|=\mu$.
\end{proof}

\section{Components of subdiagrams and their stabilizers}
\label{sec:quotient-components}

The exact fibers have equal cardinality, but choosing their bijections
independently would not preserve assignments when a diagram is replaced
by a subdiagram. We first show that restriction identifies entire
exact fibers, not merely subsets of them. Following these canonical
identifications groups the fibers into components. Only after forming
the components will we choose bijections; their stabilizers determine
which symmetries those choices must respect.

\subsection{Contractions and canonical identifications}

A vertex is a pair $v=(r,U)$, where $r\le q$, $q=r_a$, and $U$ is a
nonempty ideal of $\downarrow a$. Vertices form a ground set by
Lemma~\ref{lem:availability-ideals}. For $n\in U$ define a
\emph{contraction}
\begin{equation}\label{eq:quotient-contraction}
 (r,U)\longrightarrow(r_n,U_n),\qquad
 U_n=\{x\sqsubseteq c:nx\in U\},\quad nc=a.
\end{equation}
Left multiplication preserves conditional joins. In fact cancellation
and the common-cone description of a join give
$n(x\vee y)=nx\vee ny$ whenever these elements have a common right
multiple. It follows that $U_n$ is an ideal. Contractions compose;
there is at most one contraction between fixed endpoints, since the
center of $r_n$ is $r(n)$ and $r$ is injective. Write $B_v(F)$ and
$L_v$ for the two fibers at a vertex.

\begin{lemma}[Contraction identifications]\label{lem:contraction-fibers}
For every contraction $v\to w$, the label fibers satisfy $L_v=L_w$,
and inclusion of symbolic elements gives a bijection
$B_w(F)\longrightarrow B_v(F)$. These identifications commute with
composition.
\end{lemma}
\begin{proof}
Let $v=(r,U)$, $w=(r_n,U_n)$, and first take a symbolic class
available at $r_n$. Its full availability ideal $J$ at $r$ contains
$n$. For $z\sqsubseteq a$, ideal closure gives
\[
 z\in J\quad\Longleftrightarrow\quad z\vee n\in J.
\]
Write $z\vee n=nx$. Availability at $r_{z\vee n}$ is exactly the
corresponding availability at the subdiagram $(r_n)_x$ of $r_n$.
Thus the restricted ideal at $r_n$ determines the full ideal at $r$.
If that restricted ideal is $U_n$, the last test is equivalent to
$z\vee n\in U$, and hence to $z\in U$. Conversely, an element of
$B_v(F)$ is available at $r_n$ because $n\in U$, and its restricted
ideal is $U_n$. This proves surjectivity; injectivity follows from
Equation~\eqref{eq:symbolic-inclusion}. For labels use the same
argument with literal membership in ranges. Commutation follows
because all the symbolic maps retain the same representative names.
\end{proof}

Let a component be a connected component of the vertex system when
the direction of contractions is ignored. The successors of each
vertex are directed: contract at the join of two contraction indices.
Any two vertices in one component have a common successor. Indeed,
this holds along a single edge; along a finite path combine a common
successor of its initial part with the final edge, using directedness
of successors at the shared vertex. The same induction gives a common
successor for any finite set of vertices in a component.

Consequently all finite chains of the identifications in
Lemma~\ref{lem:contraction-fibers} are consistent. One can compare a
finite chain at a common successor, where both composites are the
same map. In a component $C$, glue the disjoint union of its $B_v(F)$
by these identifications and call the resulting set $B_C(F)$. Each
fiber maps bijectively onto $B_C(F)$. In particular, the gluing creates
no new identifications within one fiber. All label fibers in $C$
are the same ground subset of $I$; denote it by $L_C$. Both sets have
cardinality $\mu$ in $V[F]$.

Choose in $V$ a vertex $(r,U)$ of $C$ and a nondecreasing cofinal
sequence $(n_i:i<\omega)$ in $U$, with $n_0=1$. Let $v_i$ be its
contraction at $n_i$. This chain is cofinal in the entire component.
To check this, compare any vertex with $v_0$ at a common successor,
which is a contraction of $v_0$ at some $n\in U$. Choose $i$ with
$n\sqsubseteq n_i$. Then $v_i$ is a successor of that common
successor and hence of the given vertex.

\subsection{Only countably many residual symmetries}

The group $G_q$ acts on the ground vertex system. On symbolic fibers
it acts by transport between generics:
\begin{equation}\label{eq:symbolic-transport}
 T_\pi^F:\mathcal D_F(r)\longrightarrow
                   \mathcal D_{\pi F}(\pi r),
 \qquad [\tau]_r\longmapsto[\pi\tau]_{\pi r}.
\end{equation}
The forcing symmetry lemma makes this well defined, and it commutes
with all contraction identifications. We use the same notation for
transport on the glued fibers, and suppress the superscript when its
source generic is specified by the argument.

Let $G_C$ be the setwise stabilizer of $C$ in $G_q$. Define
\[
 K_C=\bigcup_{i<\omega}\Fix(D_{r_{n_i}}).
\]
Each displayed fixer is a subgroup of $G_q$, because $D_q$ is contained
in every chain range. It fixes a vertex and therefore lies in $G_C$.
The ranges decrease, so the subgroups increase and their union is a
subgroup. For $\sigma\in G_C$, the moved chain is also cofinal.
For every $i$ there is $j$ with
$D_{r_{n_j}}\subseteq D_{\sigma r_{n_i}}$. Hence
\[
 \sigma\Fix(D_{r_{n_i}})\sigma^{-1}
 =\Fix(D_{\sigma r_{n_i}})\subseteq\Fix(D_{r_{n_j}}).
\]
Applying the same reasoning to $\sigma^{-1}$ proves that $K_C$ is
normal in $G_C$.

\begin{lemma}[Countable stabilizer quotient]\label{lem:countable-germs}
In the ground model, $|G_C/K_C|\le\aleph_0$.
\end{lemma}
\begin{proof}
Choose the countable-base tower of Theorem~\ref{thm:monoid} so that
its initial submonoid $T_0$ contains all $n_i$ and all right
quotients between them. For $\pi\in G_C$, consider its alignments
with contractions of the chosen vertex:
\[
 \mathcal A_\pi=
 \{(l,m)\in\omega\times U:\pi r_{n_l}=r_m\}.
\]
This set is nonempty: the moved chain is cofinal in $C$, so one of
its terms is a successor of $v_0$. Every second coordinate $m$ of
an alignment divides some $n_N$, by cofinality of $(n_i)$ in $U$.

To apply Lemma~\ref{lem:mono-descent}, fix $(l,m)\in\mathcal A_\pi$
and any $N$. Choose $l'\ge l$ so that $\pi v_{l'}$ is a successor
of $v_N$. Its diagram is $r_k$ for some $k\in U$ with
$n_N\sqsubseteq k$. Let $d\in T_0$ be the right quotient
$n_ld=n_{l'}$. Then
\[
 r_k=\pi r_{n_{l'}}=(\pi r_{n_l})_d=r_{md}.
\]
Injectivity of $r$ gives $k=md$. Consequently $md\in U$,
$n_N\sqsubseteq md$, and $(l',md)\in\mathcal A_\pi$.
These are exactly the replacement hypotheses of
Lemma~\ref{lem:mono-descent}. Its conclusion supplies an alignment
$(l,m)$ with $m\in T_0$.

There are countably many resulting alignment pairs
$(l,m)\in\omega\times T_0$. Two permutations giving the same pair
agree pointwise on $D_{r_{n_l}}$; if they are $\pi$ and $\rho$,
then $\pi^{-1}\rho\in\Fix(D_{r_{n_l}})\subseteq K_C$.
Enumerate the pairs, and assign to each coset the first
pair realized by any of its members. Distinct cosets receive distinct
pairs. This proves the bound.
\end{proof}

\section{Equivariant assignments and the quotient bijection}
\label{sec:quotient-equivariance}

For each component $C$ we seek a bijection
$f_C(F):L_C\to B_C(F)$ satisfying
\begin{equation}\label{eq:component-covariance}
 f_C(\sigma F)(\sigma j)=T_\sigma^F(f_C(F)(j))
              \qquad(\sigma\in G_C).
\end{equation}
This is an identity between computations at different generics, not
an action on one fixed extension. We first construct a bijection
covariant under $K_C$: its elements fix a sufficiently late subdiagram,
so they fix the tails of suitable name codes. We then handle the
remaining countable quotient $G_C/K_C$. Bounded Cohen-row patterns
distinguish its cosets and select a transported bijection with the
full covariance in Equation~\eqref{eq:component-covariance}.

\subsection{An assignment invariant under a tail fixer}

For the chosen chain of $C$, let $\mathcal Q_C$ be the ground set of
sequences
\[
 c=(\tau_i:i<\omega),\qquad \tau_i\in\mathcal X(r_{n_i}).
\]
Say that $c$ represents $b\in B_C(F)$ if, from some index onward,
$\tau_i$ represents an element of the exact fiber $B_{v_i}(F)$ and
its image in $B_C(F)$ is $b$. Finite initial segments are ignored.
Every $b$ has such a code: choose representatives at all chain
vertices in $V[F]$, using their bijections with $B_C(F)$. This is a
countable sequence of ground name codes, and hence belongs to $V$
by countable closure of $\mathbb R$.

Fix a ground wellordering of $\mathcal Q_C$, of order type $\lambda_C$.
Assign to $b$ its least representing code. Distinct $b$ have different
codes, so this defines a wellorder $<_F$ of $B_C(F)$.

\begin{lemma}[Covariance of the least-code order]
\label{lem:least-code-covariance}
For every $k\in K_C$, transport $T_k^F$ is an order isomorphism
from $(B_C(F),<_F)$ to $(B_C(kF),<_{kF})$. The group $K_C$ fixes
$L_C$ pointwise.
\end{lemma}
\begin{proof}
Choose $N$ such that $k$ fixes $D_{r_{n_N}}$ pointwise. For $i\ge N$
it fixes $r_{n_i}$ and every name in $\mathcal X(r_{n_i})$ literally.
Consequently each particular ground code $c$ represents $b$ at $F$
if and only if the \emph{same code} represents $T_k^F(b)$ at $kF$.
The forcing symmetry lemma proves this on the tail; the definition
ignores the remaining finite prefix. Thus the sets of functioning
codes are identical, and their minima under the fixed ground
wellorder are identical. This proves the order assertion. Also
$L_C\subseteq D_{r_{n_i}}$ for every $i$, so $k$ fixes its every
member.
\end{proof}

Let $\theta_C(F)$ be the order type of this order. It is at most
$\lambda_C$, since it is the order induced on a subset of the ordered
code set. It has cardinality $\mu$ in $V[F]$. Preservation of $\mu$
and $\mu^+$ implies
\[
 \mu\le\theta_C(F)<\mu^+,
 \qquad |\theta_C(F)|^V=\mu.
\]
Here the order type is an ordinal, hence a ground ordinal. The two
preserved cardinal endpoints justify the last equality even if smaller
cardinals have collapsed. In $V$, choose bijections
$e_\theta:\mu\longrightarrow\theta$ for the set of ordinals
$\theta\le\lambda_C$ with ground cardinality $\mu$. Also fix a ground
bijection $l_C:\mu\longrightarrow L_C$. If
$b_F:\theta_C(F)\longrightarrow B_C(F)$ is the unique order
isomorphism, define
\[
 f_C^0(F)=b_F\circ e_{\theta_C(F)}\circ l_C^{-1}.
\]
This is a bijection $L_C\to B_C(F)$. All choices have been restricted
to sets; no simultaneous choice over a proper class of ordinals occurs.
Lemma~\ref{lem:least-code-covariance} gives
\begin{equation}\label{eq:preliminary-covariance}
 f_C^0(kF)(j)=T_k^F\bigl(f_C^0(F)(j)\bigr)
       \qquad(k\in K_C,\ j\in L_C).
\end{equation}
The ground wellorder of all codes need not be invariant under $k$.
The proof uses eventual fixation of each code's entries instead.

\subsection{Detecting and choosing a coset}

Put $\Gamma_C=G_C/K_C$, writing cosets as $\pi K_C$. It is a
countable ground
group by Lemma~\ref{lem:countable-germs}. For each nonidentity coset
$\delta$, choose a ground representative $\gamma_\delta$ and points
\[
 s_{\delta,i}\in D_{r_{n_i}},\qquad
 \gamma_\delta(s_{\delta,i})\ne s_{\delta,i}
       \quad(i<\omega).
\]
Such a point exists for every $i$, since otherwise
$\gamma_\delta\in K_C$. Fix an enumeration of the set of nonidentity
cosets. For $\alpha<\kappa$ and $\pi\in G_C$, define the pattern
\begin{equation}\label{eq:detector-pattern}
 P_F^\alpha(\pi K_C)=
 \left(
  \left[\left(A_{\pi(s_{\delta,i})}^F
                       \restriction\alpha:i<\omega\right)\right]_{\rm ev}
              :\delta\in\Gamma_C\setminus\{1\}\right),
\end{equation}
where $[\cdot]_{\rm ev}$ denotes equivalence modulo eventual equality
of sequences. We regard each row as its binary characteristic
function. The equivalence classes in this formula are taken in the
ground set of sequences of subsets of $\alpha$.

This ground interpretation is legitimate. Every $\alpha<\omega_1$
is countable. The bounded row restrictions and all their countably
many coordinates, for all detectors, form a countable sequence of
ground bits and indices. Since $\mathbb R$ adds no such sequences,
these data belong to $V$. Eventual equality between ground sequences
has the same meaning in the extension. For all $\alpha<\kappa$
together the possible patterns form a ground set, on which choose a
ground wellorder, separately for each $\alpha$ if desired.

\begin{lemma}[A covariant coset selector]\label{lem:coset-selector}
Formula~\eqref{eq:detector-pattern} is independent of the
representative of $\pi K_C$. For each $F$ there is an
$\alpha<\kappa$ at which all distinct cosets have distinct patterns.
There is a definable selection $c_C(F)\in\Gamma_C$ satisfying
\begin{equation}\label{eq:coset-selector-covariance}
 c_C(\sigma F)=\sigma c_C(F)\qquad(\sigma\in G_C).
\end{equation}
\end{lemma}
\begin{proof}
If $k\in K_C$, it fixes all $s_{\delta,i}$ once $i$ is sufficiently
large. Replacing $\pi$ with $\pi k$ therefore changes each detector
sequence in only finitely many places, proving representative
independence.

Suppose $\pi K_C\ne\pi'K_C$. Use the detector for
$\delta=(\pi^{-1}\pi')K_C$ and write
$\pi^{-1}\pi'=\gamma_\delta k$ with $k\in K_C$. Eventually
$k(s_{\delta,i})=s_{\delta,i}$, and hence
\[
 \pi(s_{\delta,i})\ne\pi'(s_{\delta,i}).
\]
Distinct label indices give distinct Cohen rows. For each of these
tail positions choose a bit where the two rows differ. The countably
many bit ordinals are bounded below regular $\kappa$. Above that
bound the two restricted sequences disagree at every position on a
tail, so are unequal modulo eventual equality. There are countably
many pairs of cosets; a further countable supremum supplies one bound
that distinguishes all pairs. These choices take place in $V[F]$.
They use regularity of $\kappa$, not closure of $\mathbb P$.

At the least distinguishing $\alpha$, choose the coset whose pattern
is least in the fixed ground wellorder. If $\Gamma_C$ is trivial,
choose its sole element and use $\alpha=0$. For $\sigma\in G_C$,
label transport gives, for every $s\in I$,
$A_{\sigma s}^{\sigma F}=A_s^F$. Thus
\begin{equation}\label{eq:pattern-covariance}
 P_{\sigma F}^\alpha((\sigma\pi)K_C)
                    =P_F^\alpha(\pi K_C).
\end{equation}
The actual ground pattern objects on the two sides are equal. Left
multiplication permutes the coset set, so the least distinguishing
bound is unchanged and the same pattern object is minimal. This
proves Equation~\eqref{eq:coset-selector-covariance}.
\end{proof}

For a coset $\gamma=\pi K_C$, define a transported preliminary
bijection by
\begin{equation}\label{eq:coset-assignment}
 f_C^\gamma(F)(j)=
 T_\pi^{\pi^{-1}F}
       \left(f_C^0(\pi^{-1}F)(\pi^{-1}j)\right),\qquad j\in L_C.
\end{equation}
Equivalently, the square in Figure~\ref{fig:moved-generic-transport}
commutes.

\begin{figure}[htbp]
\centering
\begin{tikzpicture}[font=\small,>=Stealth]
  \node (a) at (0,1.65) {$L_C$};
  \node (b) at (6,1.65) {$B_C(\pi^{-1}F)$};
  \node (c) at (0,0) {$L_C$};
  \node (d) at (6,0) {$B_C(F)$};
  \draw[->] (a) -- node[above] {$f_C^0(\pi^{-1}F)$} (b);
  \draw[->] (a) -- node[left] {$\pi\restriction L_C$} (c);
  \draw[->] (b) -- node[right] {$T_\pi^{\pi^{-1}F}$} (d);
  \draw[->] (c) -- node[below] {$f_C^{\pi K_C}(F)$} (d);
\end{tikzpicture}
\caption{Transporting a preliminary bijection, for $\pi\in G_C$.
The square commutes: to assign the bottom label $j$, first compute
at $\pi^{-1}j$ in the generic $\pi^{-1}F$, then transport the result
to $F$. The top and bottom computations use $\pi^{-1}F$ and $F$,
respectively.}
\label{fig:moved-generic-transport}
\end{figure}


This is independent of $\pi$. Replacing $\pi$ with $\pi k$ inserts
$k^{-1}$ in the source generic, and
Equation~\eqref{eq:preliminary-covariance} cancels its transport;
$k$ fixes every label of $L_C$. Define
$f_C(F)=f_C^{c_C(F)}(F)$. Equations~\eqref{eq:coset-selector-covariance}
and~\eqref{eq:coset-assignment} now give the required
Equation~\eqref{eq:component-covariance}.
This calculation explains why the preliminary computation in
Equation~\eqref{eq:coset-assignment} is made at $\pi^{-1}F$.
An arbitrary transport of the computation at $F$ would not give the
required identity.

\subsection{Transport over all components}

Choose in $V$ one representative from each $G_q$-orbit of components
and perform the preceding set-sized choices for these representatives.
For $C'=\sigma C$, transport $f_C(\sigma^{-1}F)$ by $\sigma$ to
obtain $f_{C'}(F)$. Two possible choices of $\sigma$ differ by an
element of $G_C$, so Equation~\eqref{eq:component-covariance} proves
that this definition is independent of the choice. It yields
bijections for every component, covariant under all $G_q$. Pulling
back along the canonical maps $B_v(F)\to B_C(F)$ gives bijections
\begin{equation}\label{eq:coherent-fiber-bijections}
 f_{r,U}(F):L_{r,U}\longrightarrow B_{r,U}(F).
\end{equation}
They commute with every contraction, and their transports under every
$\pi\in G_q$ are the corresponding bijections at $\pi F$.

We record the definability needed to convert these computations into
one symmetric name. The ground sets of vertices, components, name
codes, code wellorders, patterns and auxiliary bijections are all
sets. The sections of the product names can be constructed in an
$\mathbb R$-extension from its generic filter; the ordinary forcing
relation for the ground set poset $\mathbb P$ is definable there.
Symbolic quotients and their gluing use equivalence relations on sets.
All remaining steps use set operations, least elements of prescribed
wellorders, and unique order isomorphisms. They therefore define,
by an ordinary set-theoretic formula in an $\mathbb R$-extension, the
predicate
\begin{equation}\label{eq:assignment-predicate}
 \Phi(F,r,j,\tau):
 \quad \tau\in\mathcal X(r)\text{ represents }
 f_{r,U_j}(F)(j)\text{ in }\mathcal D_F(r).
\end{equation}
Here $U_j$ is ground data from Equation~\eqref{eq:label-availability}.
One formal implementation uses a ground set of all normalized name
codes and its ordinary name-evaluation relation to construct their
$\mathbb P$-sections. The variable for a ground code and the actual
forcing name that it codes are kept distinct. In particular, we never
use the false assertion that an automorphism moves a check-name for a
ground code to the check-name for its transported code.

The proved semantic covariance is the identity
\[
 \Phi(F,r,j,\tau)\quad\Longleftrightarrow\quad
 \Phi(\pi F,\pi r,\pi j,\pi\tau)
       \qquad(\pi\in G_q).
\]
These computations are ordinary forcing arguments for this fixed
formula. More explicitly, let
\[
 \dot F^{\pi}=
 \{(\check{\pi d},d):d\in\mathbb R\},
\]
whose value at $F$ is $\pi F$. The preceding covariance proof gives
the equivalence, forced by $d_0$ in $\mathbb R$, between
$\Phi(\dot F,r,j,\tau)$ and
$\Phi(\dot F^{\pi},\pi r,\pi j,\pi\tau)$, with all auxiliary
ground choices fixed. Under the name action,
$\pi\dot F^{\pi}=\dot F$ literally, while check names for the
displayed ground arguments remain fixed. Applying ordinary forcing
symmetry therefore gives
\begin{equation}\label{eq:assignment-forcing-covariance}
 d\forces_{\mathbb R}\Phi(\dot F,r,j,\tau)
 \quad\Longleftrightarrow\quad
 \pi d\forces_{\mathbb R}
                  \Phi(\dot F,\pi r,\pi j,\pi\tau)
 \qquad(d\le d_0).
\end{equation}
All auxiliary ground choices are already fixed in the defining
formula; covariance was proved for this one construction. The
$\mathbb R$-forcing values in this display are ground sets of
conditions. Equivalently, if $F'$ contains $\pi d$, put
$F=\pi^{-1}F'$ and apply covariance to the computation at $F$;
the converse uses $\pi^{-1}$. This generic description abbreviates
the fixed-formula argument above and requires no transitive-ground
assumption in the consistency deduction.

\subsection{One hereditarily symmetric graph}

\begin{theorem}[Quotient classification]\label{thm:quotient-classification}
In $W$, every nonwellorderable surjective image of $A$ is in bijection
with $A$.
\end{theorem}
\begin{proof}
We first finish the construction for the equivalence relation fixed
at the start of Section~\ref{sec:quotient-fibers}. Use canonical
equivariant ordered-pair names and form the product name
\begin{equation}\label{eq:quotient-bijection-name}
 \dot J=
 \left\{
 \left(\langle\dot A_j,\tau\rangle^{\bullet},(d,r)\right):
 \begin{array}{l}
 r\le q,\ j\in D_r,\ \tau\in\mathcal X(r),\ d\le d_0,\\
 d\forces_{\mathbb R}\Phi(\dot F,r,j,\tau)
 \end{array}\right\}.
\end{equation}
The superscript denotes the name for the ordered pair. This is a set
name: every index ranges over a ground set and the defining forcing
relation is a ground relation. Equation~\eqref{eq:assignment-forcing-covariance}
sends each entry under $\pi\in G_q$ to another entry. Applying
$\pi^{-1}$ gives equality of the sets of entries, so $G_q$ fixes
$\dot J$ literally. Each member name is a pair of HS names; thus
$\dot J$ is HS with support $D_q$. The individual member names may
have different supports. Hereditary symmetry does not require all of
them to use $D_q$.

Take a product generic $(F,H)$ containing $(d_0,q)$. For every
$r\in H$ below $q$, the graph contains exactly the local assignments
of Equation~\eqref{eq:coherent-fiber-bijections}. In detail, whenever
$\Phi(F,r,j,\tau)$ holds, the $\mathbb R$ truth lemma gives a
condition of $F$ forcing it; strengthening with $d_0$ retains a
condition of $F$, and Equation~\eqref{eq:quotient-bijection-name}
includes its entry. Conversely, every active entry forces its
asserted assignment. All classes are represented by names in their
symbolic equivalence classes, so the local assignments are complete.

We check the comparison when the diagram is strengthened. Suppose
$t\le r\le q$, $r=t_b$, and $j\in D_r$. The availability ideal
of $j$ at $t$ contains $b$. Hence
\[
 (t,U_j\text{ computed at }t)
        \longrightarrow(r,U_j\text{ computed at }r)
\]
is a contraction. Equation~\eqref{eq:coherent-fiber-bijections}
therefore gives the same assigned symbolic element under inclusion.
In particular, its product-name evaluations are equal in every
$\mathbb P$-generic containing $t$.

Any two active entries of $\dot J$ occur together below a common
stronger diagram in $H$. At that diagram the union of the
exact-fiber bijections is a bijection
$D_t\to\mathcal D_F(t)$: availability ideals partition both sets.
The preceding comparison proves single-valuedness and injectivity of
$\dot J^{F,H}$. Distinct symbolic elements are forced unequal there
by local decision, and distinct indices have distinct Cohen labels.
Thus these are also single-valuedness and injectivity for the actual
sets of the graph.

Every label index belongs to a diagram in $H$ below $q$, by the
density of extending diagrams to contain prescribed supports.
Therefore the graph has domain $A$. To reach an arbitrary $E$-class,
take any label belonging to that class and include its index in such
a diagram $r$. Its canonical supported class name represents an
element of $\mathcal D_F(r)$. The local fiber bijection is onto the
fiber containing that element, so the graph reaches the class.
This is a pointwise existence argument; no simultaneous selection
of representatives of all classes has been made. The range is
exactly $A/E$. The graph belongs to $W$ because its name is HS.

The construction was made on an extension of an arbitrary symmetric
condition asserting the nonwellorderable-quotient hypothesis. Applying
the symmetric forcing theorem and truth lemma therefore proves in
$W$ that every such quotient $A/E$ is in bijection with $A$.
Finally, for a surjection $h:A\to Y$ in $W$, form its kernel
relation $E$ in $W$. The map $[a]_E\mapsto h(a)$ is a bijection
$A/E\to Y$, defined by its unique values without choosing class
representatives. If $Y$ is nonwellorderable, so is $A/E$, and composing
this bijection with the one just constructed proves the theorem.
\end{proof}
